\subsection{Exact unit norm images in the ninth-root comparison field}

\begin{statement}{A.203}{Proved}
{Two evaluated unit-support kernels and the induced relative-unit character}
Retain the embedded fields and generators
\[
\begin{gathered}
 F=\Q(\alpha,\zeta),\quad F'=\Q(\alpha,z),\quad
 E_2=\Q(t,\zeta),\quad M=\Q(t,z),\\
 t^3=d^2\alpha,\quad \alpha^3=3,\quad d=649,
 \quad z=\zeta_9,\quad\zeta=z^3,\\
 \sigma(t)=\zeta t,\quad\sigma(z)=z,
 \qquad h(t)=t,\quad h(z)=z^4.
\end{gathered}
\tag{A.203.1}
\]
Thus \(\sigma\) generates \(\operatorname{Gal}(M/F')\),
and \(h\) generates \(\operatorname{Gal}(M/E_2)\).
Write
\(\mathcal U_L=\widehat{\mathcal O_L^\times}_3\),
\(\mathcal F_L=\widehat{L^\times}_3\), and use the actual
field norm in each of the following maps.
The completed unit norm quotients are
\[
 \mathcal U_{E_2}/N_h\mathcal U_M\simeq(\Z/3)^2,
 \qquad
 \mathcal U_{F'}/N_\sigma\mathcal U_M\simeq(\Z/3)^5.
\tag{A.203.2}
\]
Every element of \(\mathcal U_{E_2}\) is a norm from
\(\mathcal F_M\).  The first quotient in (A.203.2) has
basis the classes of the actual units \(r_3,r_4\) of
A.175.  If
\(\mathcal V=\mathcal U_{F'}\cap N_\sigma\mathcal F_M\),
then the literal inclusion of units in the multiplicative
completion gives the exact sequence
\[
 0\longrightarrow\mathcal V/N_\sigma\mathcal U_M
 \longrightarrow\mathcal U_{F'}/N_\sigma\mathcal U_M
 \longrightarrow\mathcal U_{F'}/\mathcal V
 \longrightarrow0,
\tag{A.203.3}
\]
whose three terms are respectively
\((\Z/3)^2,(\Z/3)^5,(\Z/3)^3\).
The left term has basis \(e_2,e_3/e_4\), with the A.191
unit conventions.  The same finite quotients and these
representatives occur for actual global units and actual
field norms before taking pro-three completions.

There is also an evaluated relative-unit map.  Let
\(\tau\) be complex conjugation, put
\(\mathcal U_{E_2}^-=\{a:\tau a=a^{-1}\}\), and let
\(\mathcal R_1\) be the completed global-unit kernel of
\(N_{K_2/K_1}\).  For \(v\in\mathcal R_1\), choose
\(a\in\mathcal U_{E_2}^-\) with
\(v=\sigma^2(a)/\sigma(a)\).
The class of \(a\) in the first quotient of (A.203.2)
is independent of this choice.  It defines an onto map
\[
 \eta:\mathcal R_1\longrightarrow(\Z/3)^2,
 \qquad
 \eta(v_2)=(0,2),\quad
 \eta(v_3)=(1,0),\quad
 \eta(v_4)=(0,1),
\tag{A.203.4}
\]
where
\(v_2=u_2/U^{19}\), \(v_3=u_3/U^{-1}\), and
\(v_4=u_4/U^{538}\).
These are the actual relative units of A.159 and A.191.
The map in (A.203.4) is specified by the displayed
cyclotomic unit norm quotient and Hilbert--90 difference
map.
\end{statement}

\begin{proof}
\emph{The target restriction maps.}
The complete pro-three unit bases of A.175 and A.191 give
the coordinate orders
\[
\begin{aligned}
 E_2:&\quad (u_1,u_2,u_3,u_4,\epsilon,r_2,r_3,r_4;
                  \zeta),\\
 F':&\quad (e_1,\ldots,e_9).
\end{aligned}
\tag{A.203.5}
\]
In the first row the first eight are free coordinates and
the last is its order-three torsion coordinate.
Restriction of Kummer classes from \(E_2\) to \(M\)
has kernel generated by \(\zeta\), because
\(M=E_2(\zeta^{1/3})\) is this actual cyclic Kummer
extension.  A unit whose Kummer class is in that kernel
is \(\zeta^i\) times a cube of an actual unit: its
cube root has zero valuation at every finite prime.
Consequently the restriction kernel on unit classes
modulo cubes is exactly that one-dimensional line.
The literal norm is
\[
                  N_h(z)=z^{1+4+7}=\zeta.
\tag{A.203.6}
\]
Thus this whole kernel is already in the unit norm image.

For \(F'\to M\), the Kummer kernel on field classes is
generated by \(d^2\alpha\).  At any prime above eleven,
this element has valuation two: \(F'/\Q\) is unramified
there.  A class \((d^2\alpha)^i\), \(i=1,2\), cannot
be represented by a global unit modulo cubes, since its
valuation is nonzero modulo three.  Restriction on the
nine-dimensional unit quotient is therefore injective.
These two restriction maps are exactly the target maps
used in the following finite computation.

\emph{Norms of a complete source basis.}
Use the 51 actual units in the order (A.199.6), and the
27 residue maps (A.199.7).  The latter give an isomorphism
on the entire unit quotient modulo cubes by A.199.
For a residue map \((t,z)\mapsto(a,b)\), the two norm
evaluations are the literal products
\[
\begin{aligned}
 N_h(w)(a,b)&=\prod_{j=0}^2w(a,b^{4^j}),\\
 N_\sigma(w)(a,b)&=\prod_{j=0}^2w(ab^{3j},b).
\end{aligned}
\tag{A.203.7}
\]
These formulas retain the full action on every old and
new unit, including both conjugate-root orbits.  The
certificate evaluates them on all 51 generators at all
27 maps.

After restriction, express \(N_h(w)\) in the first
eight coordinates of (A.203.5); the ninth coordinate is
precisely the kernel already handled in (A.203.6).
The columns \(w=e_5,e_7,r,s,g(s),g^2(s)\) have matrix
\[
 H=\begin{pmatrix}
 1&0&1&2&2&1\\
 0&0&1&0&1&2\\
 0&0&0&2&0&2\\
 0&0&1&0&0&2\\
 2&2&2&0&0&0\\
 0&0&0&1&0&2\\
 0&0&0&0&0&0\\
 0&0&0&1&0&2
 \end{pmatrix}\quad\text{over }\F_3.
\tag{A.203.8}
\]
It has rank six.  Every one of the 51 norm columns
satisfies the two exact equations
\[
                       x_{r_3}=0,\qquad x_{r_4}=x_{r_2}.
\tag{A.203.9}
\]
These equations cut out a dimension-six subspace, so
(A.203.8) proves equality with the entire restricted
norm image.  Adding (A.203.6) proves that its dimension
before restriction is seven, with quotient map
\[
  (x_{u_1},\ldots,x_{u_4},x_\epsilon,x_{r_2},x_{r_3},
                x_{r_4};x_\zeta)
                  \longmapsto(x_{r_3},x_{r_4}-x_{r_2}).
\tag{A.203.10}
\]

For \(N_\sigma\), use the nine \(e_i\) coordinates.
The columns \(w=r,s,g(s),g^2(s)\) are
\[
 Z=\begin{pmatrix}
 0&1&0&1\\
 2&1&2&2\\
 0&2&1&1\\
 0&1&2&0\\
 1&2&0&1\\
 2&0&2&1\\
 0&0&0&0\\
 1&1&2&0\\
 2&0&1&2
 \end{pmatrix}.
\tag{A.203.11}
\]
Its rank is four.  All 51 columns are killed by the
following rank-five matrix, which therefore has exactly
the unit norm image as its kernel:
\[
 C=\begin{pmatrix}
 1&1&2&1&1&0&0&0&0\\
 1&2&0&0&0&1&0&0&0\\
 0&0&0&0&0&0&1&0&0\\
 0&1&1&2&0&0&0&1&0\\
 2&2&1&0&0&0&0&0&1
 \end{pmatrix}.
\tag{A.203.12}
\]
In particular \(C\) is an explicit coordinate quotient
map for the second group in (A.203.2).

The source units in A.199 generate the full compact
unit module, so these computations give the images of
all units.  Also both norms contain all cubes of target
units, since the scalar inclusion followed by norm is
cubing.  The compact images are closed.  Their quotients
are therefore exactly the exponent-three quotients
computed above, proving (A.203.2).

\emph{The field-norm comparison and support kernels.}
The extension \(M/E_2\) is unramified at every finite
prime except its unique prime above three, and it is
totally imaginary.  Any global unit is a local norm at
every unramified place.  Global reciprocity for this
cyclic extension makes its remaining local character
zero at three.  The cyclic Hasse norm theorem
\cite[Chapter VIII, Theorem 3.1, printed p.~234]{MilneCFT}
therefore makes every actual \(E_2\) global unit a norm
of an element of \(M^\times\).  This proves the actual
field-norm assertion before completion.

For \(M/F'\), let \(Q\) be the first three tame rows
in A.191.  The other three lie in their span.  In the
\(e_i\) coordinate order,
\[
 Q=\begin{pmatrix}
 1&0&1&1&2&1&1&1&2\\
 1&0&1&1&1&2&2&0&2\\
 1&0&1&1&1&2&1&0&2
 \end{pmatrix}.
\tag{A.203.13}
\]
It has rank three.  A.191 proves that its kernel is
exactly the actual global-unit field-norm subgroup
modulo cubes: these are all changed-prime local norm
conditions, reciprocity supplies the unique three-prime
condition, and the cyclic Hasse norm theorem supplies
the global norm.
The literal comparison of quotient maps is
\[
 Q=LC,\qquad
 L=\begin{pmatrix}
 2&1&1&1&2\\
 1&2&2&0&2\\
 1&2&1&0&2
 \end{pmatrix},\qquad \operatorname{rank}L=3.
\tag{A.203.14}
\]
Thus (A.203.3) is, in these coordinates, precisely
\(0\to\ker L\to\F_3^5\xrightarrow{L}\F_3^3\to0\).
The exponent vectors of \(e_2\) and \(e_3/e_4\) are
\((0,1,0,0,0,0,0,0,0)\) and
\((0,0,1,2,0,0,0,0,0)\).
They are killed by \(Q\), and their images under \(C\)
are independent.  They therefore form the asserted
basis of the dimension-two support kernel.
For \(M/E_2\), (A.203.10) sends \(r_3,r_4\) to
the two standard basis vectors, proving its claimed
support-kernel representatives as well.

Here completion does not change any of these finite
quotients.  The actual unit norm image contains all
cubes and has the finite exponent-three cokernel just
computed.  The actual unit field-norm subgroup is
likewise the kernel of the displayed finite local
character map.  Its pro-three completion embeds in
\(\mathcal U_L\) with exactly that kernel.  To identify
it with \(\mathcal U_L\cap N\mathcal F_M\), the forward
inclusion follows by choosing actual norm preimages
of a finite generating set and taking their completed
powers.  The reverse inclusion follows by localization:
any completed field norm has zero in each of the finite
local norm-character quotients, so it satisfies the
same displayed unit conditions.  This proves the precise
completed intersection and both versions of the support
sequences.

\emph{The relative-unit quotient.}
The full anti-unit basis of A.175 is
\(\zeta,\epsilon,r_2,r_3,r_4\).
The actual difference map
\[
 \mathcal U_{E_2}^-\longrightarrow\mathcal R_1,
             \qquad a\longmapsto\sigma^2(a)/\sigma(a)
\tag{A.203.15}
\]
lands in the stated target: its \(\sigma\)-norm is
one, and complex conjugation fixes its value because
\(\tau\sigma\tau=\sigma^{-1}\) and \(\tau a=a^{-1}\).
It sends \(r_i\) to \(v_i\).  To verify this equality
literally, cube the image and use
\(r_i^3=\sigma(u_i)/\sigma^2(u_i)\) and
\(N(u_i)=U^{3g_i}\), where
\((g_2,g_3,g_4)=(19,-1,538)\).
The cube is \((u_i/U^{g_i})^3\); both elements are
fixed by \(\tau\), so their ratio cannot be a
nontrivial cube root of unity.
The three \(v_i\) form a pro-three basis of
\(\mathcal R_1\) by A.159, so (A.203.15) is onto.

Its kernel consists precisely of the \(\sigma\)-fixed
anti-units, namely \(\mathcal U_F^-\).
This statement passes from actual unit lattices to
their pro-three completions by flatness of \(\Z_3\).
The kernel is generated by \(\epsilon,\zeta\), and
both are in \(N_h\mathcal U_M\) by
(A.203.9)--(A.203.10) and (A.203.6).
Thus the quotient map (A.203.10) descends uniquely
through (A.203.15).  Its values on the specified
\(r_i\) are exactly (A.203.4), proving the final
assertion through the displayed maps.

The executed certificate
\path{code/two_inert_ninth_root_unit_norm_maps.py}
retains both literal norm operators on all 51 source
units, every target-coordinate comparison through the
injective A.199 detector, both rank witnesses, the
annihilator matrices, and the exact quotient factor
(A.203.14).  Its output and full matrices are retained.
\end{proof}
